% from ega0-19-fr.tex
\oldpage[0\textsubscript{IV-1}]{92}

Le diagramme
\[
  \xymatrix{
    E' \ar[r]^{\psi} & E \\
    B \ar[u]^{g} \ar[r]_{\theta} & B/\mathfrak J^2 \ar[u]_{v'}
  }
\]
où $v'$ est déduit de $v$ (18.2.8), est commutatif. D'ailleurs l'image de
$\mathfrak J$ par $\psi\circ g$ est contenue dans $Q$, donc l'image de
$\mathfrak J$ par $g$ est contenue dans $P+\mathfrak L=P$, et l'image de
$\mathfrak J^2$ par $g$ est nulle. Restreignant $g$ et $\theta$ à
$\mathfrak J$, on obtient un diagramme commutatif
\[
  \xymatrix{
    P \ar[r]^{u} & Q \\
    \mathfrak J/\mathfrak J^2 \ar[u]^{w} \ar[ur]_{v}
  }
\]
ce qui prouve que le $C$-module $\mathfrak J/\mathfrak J^2$ est projectif.
\end{env}

\begin{env}[19.5.4.2]
\label{0.19.5.4.2-fr}
Prouvons en second lieu que $a)$ entraîne que $\varphi$ est bijectif, ce qui
achèvera de montrer que $a)$ entraîne $b)$. Posons
$E_n=B/\mathfrak J^{n+1}$, et désignons par $F_n$ le quotient de
$\mathbf S_C^*(\mathfrak J/\mathfrak J^2)$ par la puissance $(n+1)$-ième de
son idéal d'augmentation. Comme $\mathfrak J/\mathfrak J^{n+1}$ est nilpotent
dans $E_n$ et que $C$ est une $A$-algèbre discrète formellement lisse,
l'automorphisme identique de $C$ se factorise en
$C\xrightarrow{f}E_n\to C=E_n/(\mathfrak J/\mathfrak J^{n+1})$ où $f$ est un
$A$-homomorphisme~; d'autre part, $\mathfrak J/\mathfrak J^2$ étant un
$C$-module projectif d'après (19.5.4.1), l'automorphisme identique de
$\mathfrak J/\mathfrak J^2$ se factorise en
$\mathfrak J/\mathfrak J^2\xrightarrow{g}\mathfrak J/\mathfrak J^{n+1}
\to\mathfrak J/\mathfrak J^2$, où $g$ est une application $C$-linéaire~; à
partir de $f$ et $g$ on obtient canoniquement un homomorphisme de $C$-algèbres
$\mathbf S_C^*(\mathfrak J/\mathfrak J^2)\to E_n$, qui d'ailleurs (par
définition de $g$) s'annule sur la puissance $(n+1)$-ième de l'idéal
d'augmentation de $\mathbf S_C^*(\mathfrak J/\mathfrak J^2)$, d'où, en
passant au quotient, un $A$-homomorphisme surjectif d'algèbres
$v:F_n\to E_n$ tel que $\gr^0(v)$ et $\gr^1(v)$ soient les automorphismes
identiques de $C$ et de $\mathfrak J/\mathfrak J^2$. Par définition de
l'homomorphisme canonique $\varphi$, on voit que $\gr^j(v)=\varphi_j$ pour tout
$j\leq n$. Notons maintenant que le noyau $\mathfrak R$ de $v$ est un idéal
nilpotent de $F_n$, de sorte que $E_n$ s'identifie à $F_n/\mathfrak R$. Comme
$B$ est une $A$-algèbre formellement lisse, le $A$-homomorphisme canonique
$p_n:B\to E_n=B/\mathfrak J^{n+1}$, qui est continu, se factorise en
$B\xrightarrow{w}F_n\xrightarrow{v}E_n$, où $w$ est un $A$-homomorphisme
continu~; comme $\gr^0(v)$ est l'identité, $w(\mathfrak J)$ est contenu dans
l'idéal d'augmentation de $F_n$, d'où $w(\mathfrak J^{n+1})=0$, si bien que
$w$ se factorise en
$B\xrightarrow{p_n}B/\mathfrak J^n=E_n\xrightarrow{w'}F_n$, et le composé
$E_n\xrightarrow{w'}F_n\xrightarrow{v}E_n$ est l'identité. En outre, comme
$\gr^0(v)$ et $\gr^1(v)$ sont les automorphismes identiques de $C$ et de
$\mathfrak J/\mathfrak J^2$, il en est de même de $\gr^0(w')$ et de
$\gr^1(w')$. Comme $\gr^*(E_n)$ est engendré par $\gr^1(E_n)$,
l'homomorphisme composé
\[
  \gr^j(F_n)\xrightarrow{\varphi_j}\gr^j(E_n)
  \xrightarrow{\gr^j(w')}\gr^j(F_n)
\]
est l'identité pour tout $j\leq n$, puisque c'est vrai pour $j=0$ et $j=1$~;
faisant $j=n$, on prouve ainsi que $\varphi_n$ est injectif, ce qui achève de
montrer que $a)$ entraîne $b)$.
\end{env}

\begin{env}[19.5.4.3]
\label{0.19.5.4.3-fr}
